% Narrative of the results. Every number is a macro generated from the artifacts
% (paper/generated/numbers.tex) or lives in a generated table.

\newcommand{\ResultsHypotheses}{%
\textbf{H2 (calibration-marginal EIG vs.\ plug-in Fisher CAT)} is
\textbf{\DecHTwo}: at $N=8$ EOST-M lowered the integrated holdout log loss in
\PosHTwo\ of \NBHTwo\ benchmarks (median improvement \MedImpHTwo\%, range
\MinImpHTwo\% to \MaxImpHTwo\%), Holm-significantly in \FavHTwo\ and never
significantly in favour of the baseline. At $N=4$ the median
improvement was \MedImpHTwoNfour\% (\FavHTwoNfour/\NBHTwoNfour\ significant),
at $N=64$ only \MedImpHTwoNsixtyfour\%.
\textbf{H3 (EOST-M vs.\ naive subsampling)} is \textbf{\DecHThree}: the
integrated score error fell by a median of \MedImpHThree\% (\FavHThree\ of
\NBHThree\ benchmarks significant).
\textbf{H4 (interval calibration)} is \textbf{\DecHFour}: EOST-M reduced the
miscalibration $|\mathrm{coverage}-0.9|$ of plug-in Fisher CAT by a median of
\MedImpHFour\% (\FavHFour\ of \NBHFour\ significant).
\textbf{H1 (information vs.\ entropy)} is \textbf{\DecHOne}: under the same
calibration-marginal model, EIG lowered held-out log loss in \PosHOne\ of
\NBHOne\ benchmarks, but by a median of only \MedImpHOne\% (range
\MinImpHOne\% to \MaxImpHOne\%), Holm-significantly in \FavHOne, short of
the frozen threshold. At the other panel sizes the same contrast was
significant more often (\FavHOneNfour/\NBHOneNfour\ benchmarks at $N=4$,
\FavHOneNsixtyfour/\NBHOneNsixtyfour\ at $N=64$; Table~\ref{tab:allN}), never in
favour of entropy, but always with median gains of a few per cent. The sign
predicted by the EOST thesis is therefore observed throughout, but its size on
held-out prediction is small once calibration uncertainty is integrated
(Corollary~\ref{cor:blind}); on score estimation the gap is larger
(Table~\ref{tab:factorial}).
GSM8K is the one benchmark on which the H2--H4 contrasts, while pointing the
same way, did not reach significance. It is also the confirmatory bank on which
an eight-system panel leaves the fewest items unlocated (\GsmFracExtremeNeight\%
of items with identical panel verdicts, against a median of
\OtherFracExtremeNeight\% elsewhere), because its leaderboard population spans
the whole accuracy range, from format failures (\GsmFloorFrac\% of the models
score below 5\%) to strong models; by the mechanism of
Figure~\ref{fig:mechanism}, there is less for marginalization to correct. Math
word problems were also the heterogeneous case of the pilot program
(GSM-Hard, Section~\ref{sec:pilot}).}

\newcommand{\ResultsFactorial}{%
Because EOST-M differs from plug-in Fisher CAT in two respects---the criterion
and the measurement model---Table~\ref{tab:factorial} crosses them. At $N=8$,
replacing plug-in parameters by expected response functions improves held-out
log loss by a median of \FactEntMgLL\% even with entropy selection, whereas
switching the criterion under plug-in parameters changes little
(entropy \FactEntPiLL\%, EIG \FactEigPiLL\%). Combining both gives
\FactEigMgLL\%. EIG beat entropy on log loss in \EigVsEntMgPos\ of
\NConfBanks\ benchmarks under the marginal model and in \EigVsEntPiPos\ under
plug-in parameters. For the score, the criterion matters more: entropy
selection is the worst adaptive rule under both models (score IAE
\FactEntMgIAE\% relative to MFI:pi under the marginal model), while EIG:mg
improves it by \FactEigMgIAE\%. Under the marginal model, the four
information-based criteria (MFI, PWFI, EIG and EVR) are practically
indistinguishable: their mean held-out log losses differ by at most
\InfoSpreadLL\% and their score errors by at most \InfoSpreadIAE\%
(Table~\ref{tab:policies}), as the first-order equivalence of
Proposition~\ref{prop:local} predicts. Most of the gain over current practice
thus comes from integrating calibration uncertainty; selecting by information
rather than by entropy is what turns that model into an efficient test, and
EIG is the principled member of an equivalence class rather than a uniquely
better rule.}

\newcommand{\ResultsPanel}{%
The gain of EOST-M over plug-in Fisher CAT decreases monotonically with the
panel size on every benchmark (Figure~\ref{fig:heatmap}a): it is largest with
four calibrators and nearly vanishes with 64, the regime predicted by
Proposition~\ref{prop:cap}, since the posterior uncertainty of item difficulties
shrinks as calibrators are added. Gains in score error (panel b) are larger
but noisier across panel sizes, and the advantage of EIG over entropy
(panel c) shows no trend with $N$.}

\newcommand{\ResultsCurves}{%
Figure~\ref{fig:curves-err} shows the score error as a function of the number
of administered items. Every model-based policy starts below the naive
estimator, because the calibrated model predicts the bank before any item is
asked: after one item, the model-based error is \HeadStartOtherMin--\HeadStartOtherMax\%
of the naive error, except on GSM8K (\HeadStartGsm\%). Two regimes follow. In
the first items, plug-in policies, which take their item parameters at face
value, reduce the error fastest: at $t=5$ a plug-in policy is ahead of EOST-M
on \EarlyPluginBetter\ of \NCurveBanks\ benchmarks. Their error then levels
off. Between $t=25$ and $t=100$ the error of plug-in Fisher CAT falls by a
median of \PlateauPi\%, that of EOST-M by \PlateauMg\%, and after 100 items
EOST-M is below plug-in Fisher CAT on \FinalEigBelowMfi\ of \NCurveBanks\
benchmarks. This is the signature of a bias floor: with eight calibrators, the
predicted outcomes of the unobserved items inherit the errors of the plug-in
parameters, and administering further items does not remove them.}

\newcommand{\ResultsLate}{%
Late in the test, entropy selection under the marginal model catches up: after
100 items it has the lowest error of all policies on \EntLateBest\ of
\NCurveBanks\ benchmarks (\EntLateBestList). Figure~\ref{fig:crossover}
locates the crossover, the item after which entropy selection stays below
EIG:mg. At $N=8$ it occurs on \CrossListNeight, and not within 100 items on
\CrossNoneNeight. Across all panel sizes the crossover tends to come later the
larger the bank (Spearman $\rho=\CrossRho$ over \CrossPoints\
benchmark--panel combinations, no crossover counted as 101): it occurred within
the budget in \CrossShareSmall\% of the combinations on the four smallest banks
(500--1,172 items) and in \CrossShareLarge\% on the four larger ones. This is
the regime change described by Proposition~\ref{prop:evr}: once competence is
pinned down relative to the number of unobserved items, an item's own outcome
uncertainty dominates the expected reduction of score variance, and the larger
the bank, the later this happens. The score-targeted criterion, which
implements this trade-off under the fitted model, ended below EIG:mg on
\EvrBelowEigEntBest\ of the \EntLateBest\ benchmarks on which entropy finished
best and on \EvrBelowEigEntNot\ of the others: it moved in the predicted
direction, but only part of the way. These late-test comparisons were not
part of the frozen plan and are exploratory. Bank size is confounded with source and
domain, and \CrossNeverList\ show no crossover at any panel size, so size is at
most one determinant of when the regime changes.}

\newcommand{\ResultsEfficiency}{%
Table~\ref{tab:efficiency} converts the curves into test lengths: at $N=8$,
EOST-M matched the score error of 100 uniformly sampled items after a median of
\EffItemsMedianNeight\ items (between \EffItemsMinNeight\ and \EffItemsMaxNeight\
across benchmarks), and with the largest panels after a median of
\EffItemsMedianNmax\ items. Table~\ref{tab:policies} compares all twelve
policies: averaged over benchmarks at $N=8$, EOST-M has a score IAE of
\IAEeigmgNeight\ versus \IAEmfipiNeight\ for plug-in Fisher CAT and
\IAEnaiveNeight\ for naive subsampling (a \IAEredNaive\% reduction), and
reaches a sustained error below 0.03 after \TtolEigmg\ items on average versus
\TtolNaive\ for naive subsampling. Adaptivity itself contributes: under the
same calibration-marginal model, random item selection---a model-assisted
estimator in the spirit of IRT-based benchmark compression
\citep{polo2024tinybenchmarks}---has a score IAE of \IAErandmgNeight, and
EOST-M's error is significantly lower on \RandIaeFav\ of \RandIaeNB\
benchmarks (secondary contrast).}

\newcommand{\ResultsRanking}{%
Leaderboards need the order of systems rather than one score. In an
exploratory analysis with eight calibrators, Kendall's $\tau$ between the
score estimated after 25 items and the true full-bank accuracy, across the
targets of a benchmark, had a median of \TauEigTwentyfive\ for EOST-M,
\TauMfiTwentyfive\ for plug-in Fisher CAT and \TauNaiveTwentyfive\ for 25
uniformly sampled items (Table~\ref{tab:ranking}). Short adaptive
tests therefore rank systems far better than random subsets of the same size,
and the two adaptive tests are close to each other.}

\newcommand{\ResultsCoverage}{%
Plug-in posteriors are overconfident (Figure~\ref{fig:coverage}): at $N=8$ the
90\% interval of plug-in Fisher CAT covered the true score \CovMfipi\ of the
time on average, against \CovEigmg\ for EOST-M and \CovNaive\ for the
model-free Wilson interval of naive subsampling. Marginalization therefore
removes a large part of the overconfidence, most at small $N$, but no
model-based adaptive policy reaches nominal coverage on real data. On the
synthetic banks (Figure~\ref{fig:coverage}b), EOST-M's score coverage at $N=8$
is \SynTwoplEigMgCovNeight\ when the model is well specified,
\SynGuessEigMgCovNeight\ with guessing and \SynTwodimEigMgCovNeight\ with a
second ability dimension (plug-in Fisher CAT: \SynTwoplMfiPiCovNeight,
\SynGuessMfiPiCovNeight\ and \SynTwodimMfiPiCovNeight). Coverage thus degrades
with misspecification, and the residual gap on real data is plausibly due to
model misspecification rather than to the marginalization itself. The two
panels of Figure~\ref{fig:coverage} also move in opposite directions as the
panel grows: on the synthetic banks the coverage of the model-based intervals
rises towards the nominal level, on real data it falls. More calibrators
sharpen the item posteriors, which helps when the model is adequate and exposes
its misfit when it is not.}

\newcommand{\ResultsDiagnostics}{%
With four calibrators, \FracExtremeNfour\% of the items of a confirmatory bank
were answered identically by every calibrator, against
\FracExtremeNsixtyfour\% with 64 (Figure~\ref{fig:diag}a); the difficulty of
these items is located only by the prior. Under the empirical-Bayes prior,
plug-in discriminations above 4 are practically absent with small panels
(\FracAgtFourNfour\% of items at $N=4$) and appear only with large ones
(\FracAgtFourNsixtyfour\% at $N=64$; panel c), so the overconfidence of
plug-in testing at small $N$ comes mainly from treating poorly located
difficulties as known rather than from extreme discriminations: the expected
response function shrinks the median effective discrimination from \AmapNfour\
to \AeffNfour\ at $N=4$ (panel b), as Proposition~\ref{prop:cap} prescribes.
Across benchmarks and panel sizes, the log-loss gain of EOST-M over plug-in
Fisher CAT increases with the fraction of unlocated items (Spearman
$\rho=\MechRho$; Figure~\ref{fig:mechanism}).}

\newcommand{\ResultsSynthetic}{%
On the synthetic banks the true abilities are known (Figure~\ref{fig:synth}).
At the primary panel size $N=8$, EOST-M had the smallest mean absolute error of
the ability estimate among the four policies of the figure on all three banks
(\SynTwoplEigMgErrNeight, \SynGuessEigMgErrNeight\ and \SynTwodimEigMgErrNeight\
on the well-specified, guessing and two-dimensional banks, against
\SynTwoplMfiPiErrNeight, \SynGuessMfiPiErrNeight\ and \SynTwodimMfiPiErrNeight\
for plug-in Fisher CAT), and its 90\% credible interval covered the true
ability more often (\SynTwoplEigMgThCovNeight, \SynGuessEigMgThCovNeight\ and
\SynTwodimEigMgThCovNeight, against \SynTwoplMfiPiThCovNeight,
\SynGuessMfiPiThCovNeight\ and \SynTwodimMfiPiThCovNeight). With four
calibrators the ranking depends on the bank; from 16 calibrators on, EIG under
the marginal model stays ahead of both plug-in policies on every bank. Coverage
of the true ability falls for every adaptive policy as the panel grows: the
posterior sharpens while the calibrated ability scale, which is fixed by the
panel and the prior, only approximates the generating one. This identification
issue is why our endpoints are defined on observable quantities
(Section~\ref{sec:framework}).}

\newcommand{\ResultsSensitivity}{%
Table~\ref{tab:sens} repeats the four contrasts under the sensitivity
variants specified in the frozen plan. Removing near-duplicate calibrators (Cohen's $\kappa\ge0.8$ with the
target, on the \SensDedupNBanks\ leaderboard benchmarks) leaves the H2 gain at
$N=8$ essentially unchanged (\SensDedupHTwoNeight\%, against
\SensMainDedupBanksHTwoNeight\% for the main analysis on the same benchmarks).
The empirical-Bayes hyperprior, by contrast, is an essential part of the
method: with the fixed weakly-informative item prior, the H2 gain at $N=8$
falls to a median of \SensFixedpriorHTwoNeight\% (main: \SensMainHTwoNeight\%)
and EOST-M improves on plug-in Fisher CAT on only
\SensFixedpriorHTwoNeightPos\ of \SensFixedpriorHTwoNeightB\ benchmarks. For
items that the panel cannot locate, the expected response function is
essentially the prior predictive, so integrating over the item posterior is
only as good as the prior it integrates; this is the failure that development
had revealed (Section~\ref{sec:dev}).
\par
The one-parameter variant specified in the frozen plan turned out to be uninformative. Its
likelihood is invariant to a joint rescaling of the abilities, the
difficulties and the common discrimination; with Gaussian priors on the former
and a single log-normal prior term on the latter, the joint posterior mode
shrinks the ability scale and inflates the common discrimination (median
\OneplMedianA, up to \OneplMaxA, across calibrations). The resulting nearly
deterministic response curves make every model-based policy grossly
overconfident (EOST-M coverage \SensCovOneplNeight\ at $N=8$), so we report
this variant but do not interpret its contrasts. The two-parameter model is
protected from this degeneracy because every item carries its own prior term on
$\log a$. \ResultsOneplReg}

\newcommand{\IntroFindings}{%
In study~1 (\NSystemsTotal\ systems, leave-family-out replay), EOST-M with eight
calibrators predicted held-out verdicts better than plug-in Fisher CAT on all \NBHTwo\
benchmarks (median \MedImpHTwo\%, Holm-significant on \FavHTwo), estimated the benchmark score with a median
\MedImpHThree\% smaller error than uniform subsampling and a median \RandIaeMed\% smaller than
random items scored by the same model, and gave less overconfident intervals; the gains shrink as
the panel grows. With strict family clusters as units the prediction gain keeps its size
(\CluHTwoStrMed\%) but is significant on only \CluHTwoStrFav\ benchmarks, below the frozen
threshold. Study~2, with one target per cluster, supported the replications of all three
study-1 results and the superiority of EOST-M over the regularized Rasch-type test;
\IntroFBFinding\ in the eight-calibrator regime EOST-M was ahead of ATLAS-style CAT, tinyBenchmarks and MINCE on every benchmark,
significantly on four or five of eight. Of its \STwoNHyp\ frozen decisions,
\STwoNSupported\ were supported, and four of the unsupported superiority hypotheses missed
the threshold by one benchmark; in a secondary random-effects meta-analysis that treats
the benchmarks as independent their pooled effects exclude zero, but only that of
\MetaSrcUnsupKeptList\ does so when the five benchmarks that share their systems count as one
source. The two prediction results (against plug-in Fisher CAT and the
Rasch-type test) did not depend on how the family clusters were built. Recalibration raised the coverage of EOST-M's 90\% intervals from
\STwoCoveigmg\ to \STwoCoveigmgRR, short of the frozen criterion and without a better
interval score. A temporal validation on later SWE-bench submissions favoured EOST-M in
\TempNSig\ of \TempNTests\ contrasts.}

\newcommand{\DiscussionSummary}{%
The evidence supports the three messages of the introduction, each with a qualification.
(1) Small panels break plug-in adaptive testing: item parameters estimated from four to eight
systems and then treated as known---mainly poorly located difficulties, and inflated
discriminations (Section~\ref{sec:study1})---make Fisher-information CAT predict worse and
state intervals that cover the score only \STwoCovmfipi\ of the time; the effect replicated with independent
targets and on later systems, and it fades as the panel grows. (2) Accounting for
calibration uncertainty fixes much of this, in more than one way: integrating the item
posterior (EOST-M) and constraining the discriminations (a Rasch-type test) are two
remedies, integration predicting slightly better and constraint giving wider, better
calibrated intervals; \DiscussionFBShort\ (3) Once uncertainty is handled, the choice among
information criteria matters little (descriptively), and the original hypothesis that
information beats entropy on held-out prediction is consistent in sign but small. What remains unsolved is calibrated uncertainty: no remedy we
tried gave nominal coverage with intervals as narrow as EOST-M's.}

\newcommand{\ConclusionText}{%
Adaptive testing can make LLM evaluation cheaper, but with the small calibration panels
of new, private or expensive benchmarks it works only if the uncertainty of the item
calibration is taken seriously. Integrating that uncertainty through expected response
functions with an empirical-Bayes prior is simple and exact given the calibrators. In two
studies with frozen analysis plans, it predicted held-out verdicts better than plug-in
Fisher-information testing, and in the second study also better than a regularized
Rasch-type test, whose score errors were similar; it estimated benchmark scores more
accurately than uniform subsampling and MINCE; and, with eight calibrators, it was ahead of
random items scored by the same model, ATLAS-style adaptive testing and tinyBenchmarks on
every benchmark without meeting the frozen criterion for these three comparisons. \ConclusionFB\ The original thesis
of this project, that information rather than entropy is what an adaptive test should
maximize, holds in sign but is modest on prediction once uncertainty is integrated, and the
choice among information criteria hardly matters. The open problem is honest uncertainty:
model-based intervals remain too narrow, and recalibration from the calibration panel
improved coverage only at the price of width. Every number in this paper can be regenerated
from pinned public data with a Docker-only pipeline and frozen analysis plans, and its
qualitative statements about results are checked automatically against the artifacts.}

\newcommand{\ResultsOneplReg}{%
A post hoc variant with a strongly informative prior on the common
discrimination, $\log a\sim\mathcal N(0,0.05^2)$, is well behaved (median
common discrimination \OneplregMedianA; EOST-M coverage
\SensCovOneplregNeight\ at $N=8$). It provides the check of
Corollary~\ref{cor:blind} that the planned variant could not: when all
items share one discrimination, the median log-loss difference between EIG and
entropy at $N=8$ is \SensOneplregHOneNeight\% under the marginal model and
\SensOneplregHOnePiNeight\% under plug-in parameters, against
\SensMainHOneNeight\% and \SensMainHOnePiNeight\% with item-specific
discriminations. The information advantage of EIG over entropy therefore comes
from the heterogeneity of discriminations, as the theory says. Under this
constrained model the gain of marginalization over plug-in parameters is also
small (\SensOneplregHTwoNeight\%), an indication that much of the H2 gain
concerns discriminations that a small panel cannot estimate well;
Section~\ref{sec:strong} tests this directly.}

\newcommand{\ResultsStrong}{%
The regularized one-parameter variant suggests a competing explanation for H2:
EOST-M may beat plug-in Fisher CAT mainly because item-specific
discriminations estimated from a small panel are unreliable, a failure that a
constrained plug-in model could avoid without any marginalization. We tested
this post hoc against strong baselines on the confirmatory design (same
targets, replicates, holdouts and panels; $N\in\{4,8,16\}$), under decision
rules recorded before the runs (Table~\ref{tab:posthoc}). The recomputed
EOST-M and MFI:pi rows reproduced the confirmatory artifacts exactly
(\PhReproRows\ rows, recomputed days later after several machine restarts).
\par
Better point estimates do not fix plug-in testing: with posterior-mean
instead of posterior-mode item parameters, held-out log loss is
\PhMfieapLLNeightAbs\% worse than EOST-M's at $N=8$ (\PhMfieapLLNeightFav\ of
\PhMfieapLLNeightB\ benchmarks significant), because the posterior mean of a
poorly located item's discrimination inherits the long right tail of its
posterior. A regularized Rasch-type plug-in test---a common discrimination with
prior $\log a\sim\mathcal N(0,0.05^2)$---is a much stronger baseline. EOST-M's
held-out log loss is lower on all \PhMfioneprLLNeightB\ benchmarks, but only
by a median of \PhMfioneprLLNeightAbs\% at $N=8$ (significant on
\PhMfioneprLLNeightFav), its score error is lower on
\PhMfioneprIAENeightPos\ of \PhMfioneprIAENeightB\ (by a median of
\PhMfioneprIAENeightAbs\%, and of \PhMfioneprIAENsixteenAbs\% at $N=16$, where a single
discrimination becomes a visible misspecification), and the constrained
model's intervals are better calibrated (coverage \PhCovMfionepr\ against
\PhCovEigmg). By the rule fixed in advance, the advantage of EOST-M over
plug-in CAT is therefore largely an advantage over \emph{unregularized
item-specific discriminations}: integrating their uncertainty and constraining
them are two remedies for the same small-panel failure, with integration ahead
on prediction and score estimation, and the constrained model ahead on
interval calibration. Finally, selecting by the posterior-expected Fisher
information under the marginal model gives the same tests as EIG in practice
(median difference \PhEfiLLNeightAbs\% at $N=8$, no benchmark significant): the optimism that
Corollary~\ref{cor:jensen} attributes to averaged information did not change
which items were worth asking in our banks.
\par
Can the two remedies be combined? In a further exploratory experiment, also
planned before running, we kept EOST-M's calibration but fixed the prior
standard deviation of $\log a$ to values between 0.05 and 0.5 instead of its
empirical-Bayes estimate (median 0.68), shrinking the discriminations towards
a common value (Figure~\ref{fig:shrink}). No setting met the
criterion fixed in advance (coverage of at least 0.80 with a median held-out
log loss at most 0.5\% above EOST-M's): coverage stayed between \ShrMgCovMin\ and
\ShrMgCovMax\ (EOST-M: \ShrEostCov) while held-out log loss rose by
\ShrMgLLMin\% to \ShrMgLLMax\%. The calibration advantage of the Rasch-type
test therefore does not come from equalizing the discriminations as such.
\par
Where does it come from? A decomposition of the post hoc runs, again planned
before computing (Table~\ref{tab:covdecomp}), answers this by
the rule fixed in advance: the Rasch-type intervals are wider than EOST-M's on
all \PhThreeWiderBanks\ benchmarks (median \PhThreeWidthPct\% wider,
Holm-significant on \PhThreeWiderSig), while its score error is not lower on
\PhThreeIaeNotLower\ of the \PhThreeCovHigher\ benchmarks where its coverage is
higher. Across all benchmark--model cells, coverage is almost a monotone
function of the ratio of mean interval width to score IAE (Spearman
$\rho=\PhThreeRho$ over \PhThreeCells\ cells), and that ratio stays below the
\PhThreeGauss\ of a calibrated 90\% interval for unbiased Gaussian errors in
every cell (at most \PhThreeMaxRatio; on average \PhThreeRatioPi\ for plug-in
2PL, \PhThreeRatioEost\ for EOST-M, \PhThreeRatioRasch\ for the Rasch-type test
and \PhThreeRatioOnePrMg\ for the marginal one-parameter model). The
constrained model is better calibrated because it is less confident, not
because it is more accurate, and under adaptive selection every model we tried
states intervals too narrow for its errors. The same marginal model with
uniformly sampled items covers \CovRandmg\ of the time on the same banks, and
EOST-M covers \SynTwoplEigMgCovNeight\ on the well-specified synthetic bank
(Figure~\ref{fig:coverage}): the shortfall arises when adaptive selection meets
a misspecified model, not from either alone.}
